\begin{longtable}{|l|X|}
    \caption{Verbs} \\
    \hline
    \textbf{Verb} & \textbf{Contract} \\
    \hline
    \endfirsthead
    \multicolumn{2}{c}%
    {\tablename\ \thetable\ -- \textit{Continued from previous page}} \\
    \hline
    \textbf{Verb} & \textbf{Contract} \\
    \hline
    \endhead
    \multicolumn{2}{r}{\textit{Continued on next page}} \\
    \endfoot
    \endlastfoot
%------------------------------------------------------------------------------

\makecell{\lstinline|abort( <arg> )| \\ \lstinline|abort…()| \\ \lstinline|abort()|\pslabel{verb:abort}} & 
A method \lstinline|abort()| or one with a name prefixed with \bsq{abort} ungracefully terminates something, most probably a thread or a process determined by the arguments, and usually as a result of an error condition of some kind. In case the method does not take an argument, either the remaining part of the name determines the subject, or the owning object represents that subject. 
Also refer to \hyperref[verb:cancel]{\bsq{cancel}}, \hyperref[verb:exit]{\bsq{exit}}, \hyperref[verb:shutdown]{\bsq{shutdown}} and \hyperref[verb:terminate]{\bsq{terminate}}. \\
\hline

%------------------------------------------------------------------------------

\makecell{\lstinline|acquire( <arg> )| \\ \lstinline|acquire…()| \\ \lstinline|acquire…( <arg> )| \\ \lstinline|acquire()|\pslabel{verb:acquire}} & 
This name or prefix for the name of a method is used for methods that requests an exclusive access to something, either the owning object (then usually there is no argument given) or the object that is identified by a given argument.
    
Usually, the return value (if any) is some kind of handle that has to be released later, by a call to a \hyperref[verb:release]{\bsq{release}} method.

\index{java.util.concurrent!Semaphore}\index{java.util.concurrent!Semaphore!acquire()}\index{java.util.concurrent!Semaphore!release()}
\begin{lstlisting}[xleftmargin=.7cm]
private final Semaphore m_Semaphore = new Semaphore( 5 );
…
m_Semaphore.acquire();
try
{
    // Do something here …
}
finally
{
    m_Semaphore.release();
}
\end{lstlisting}
For more details about the class \lstinline|Semaphore| from the package \lstinline|java.util.concurrent| refer to \autocite{ORACLE_DOC_SEMAPHORE_CLASS,ORACLE_DOC_SEMAPHORE:aquire,ORACLE_DOC_SEMAPHORE:release}. \\
\hline

%------------------------------------------------------------------------------

\makecell{\lstinline|action…( <arg> )| \\ \lstinline|action…()|\pslabel{verb:action}} & 
An \bsq{action} method performs an action, usually specified by the remainder of the name or the arguments for the call. To some extent, \bsq{action} is synonym to \hyperref[verb:perform]{\bsq{perform}}\footnote{Interstingly the \Java Runtime Libary defines a lot of methods named \lstinline|actionPerformed()|.}. \\
\hline

%------------------------------------------------------------------------------

\makecell{\lstinline|add…( <arg> )| \\ \lstinline|add( <arg> )|\pslabel{verb:add}} & 
A method with this name will add the given argument value to an internal list in the object, keeping the argument as it is, so that it could be retrieved again later.
    
Usually, an \bsq{add} method would not return something, but that is not mandatory.

Refer to chapter \tqfullref{sec:NamesForCollectionMutators} where \bsq{add} methods are already mentioned. \\
\hline

%------------------------------------------------------------------------------

\makecell{\lstinline|alert…( <arg> )| \\ \lstinline|alert( <arg> )|\pslabel{verb:alert}} & 
\bsq{Alert} methods will send a message about an irregular condition to somewhere. This could be an internal monitoring facility or a human administrator that receives an email. \\
\hline

%------------------------------------------------------------------------------

\makecell{\lstinline|alter…( <arg> )| \\ \lstinline|alter( <arg> )|\pslabel{verb:alter}} & 
A method with a name prefixed by \bsq{alter} modifies something (a picture, a file, a document), but usually not the owning object. But the state the owning object determines the kind of change that is performed to the subject defined by the arguments. 

This is somehow related to \hyperref[verb:change]{\bsq{change}}. \\
\hline

%------------------------------------------------------------------------------

\makecell{\lstinline|append( <arg> )| \\ \lstinline|append…( <arg> )|\pslabel{verb:append}} & 
The name \bsq{append} for a method indicates that it will append the given argument to an extensible data structure, usually with interpreting the argument in some way. In most cases, it is not possible to retrieve the argument from the data structure afterwards. This a typical method for a builder implementation, so usually the method returns the owning object itself (\lstinline|this|).
    
A sample for this is \lstinline|java.lang.StringBuffer::append|\index{java.lang!StringBuffer!append()} \autocite{ORACLE_DOC_STRINGBUFFER_CLASS}. \\
\hline

%------------------------------------------------------------------------------

\makecell{\lstinline|build()|\pslabel{verb:build}} & 
The terminal method for a builder has this name. It returns the built object. Usually it can be called only once on the owning object.
    
An example would be \lstinline|Stream.Builder::build|\index{java.util.stream!Stream.Builder!build()} \autocite{ORACLE_DOC_STREAM_BUILDER:build}. \\
\hline

%------------------------------------------------------------------------------

\makecell{\lstinline|call()|\pslabel{verb:call}} & 
This is the name of the main method of a thread that returns a result (an implementation of the \lstinline|Callable| interface)\index{java.util.concurrent!Callable} \autocite{ORACLE_DOC_CALLABLE_INTERFACE} and should usually only used in this or a similar context.

For details on this, refer to \autocite{ORACLE_DOC_UTIL_CONCURRENT_PACKAGE}. \\
\hline

%------------------------------------------------------------------------------

\makecell{\lstinline|cancel()| \\ \lstinline|cancel…()| \\ \lstinline|cancel…( <arg> )| \\ \lstinline|cancel( <arg> )|\pslabel{verb:cancel}} & 
A \bsq{cancel} method is used to terminate a process or operation prematurely, but not necessary because of an error condition (this would be a \hyperref[verb:abort]{\bsq{abort}} method).

The respective operation is either determined by the arguments, or it is represented by the owning object.

\bsq{Cancel} is also often used in the context of user interfaces.  

Also refer to \hyperref[verb:exit]{\bsq{exit}}, \hyperref[verb:shutdown]{\bsq{shutdown}} and \hyperref[verb:terminate]{\bsq{terminate}}. \\
\hline

%------------------------------------------------------------------------------

\makecell{\lstinline|change…( <arg> )| \\ \lstinline|change( <arg> )|\pslabel{verb:change}} & 
A \bsq{change} method implies a non-permanent modification to something. This distinguishes it from an \hyperref[verb:alter]{\bsq{alter}} method. 

The subject for the change is usually determined by the arguments, but it could be also the owning object. If given, the remainder of method's name tells the kind of change. \\
\hline

%------------------------------------------------------------------------------

\makecell{\lstinline|check()| \\ \lstinline|check( <arg> )| \\ \lstinline|check…()|\pslabel{verb:check}} & 
Usually \bdq{to check} is understood as testing for a given condition, but sometimes it is also used if an object should be marked (as used in \bdq{CheckBox}). I do not recommend the latter usage (use \hyperref[verb:mark]{\bsq{mark}} instead), and because of the chance to generally misinterpret the verb, I do not recommend to use \bsq{check} at all.

If used anyway, the respective method may not modify the object, and it should return a boolean. \\
\hline

%------------------------------------------------------------------------------

\makecell{\lstinline|clear()| \\ \lstinline|clear…()| \\ \lstinline|clear…( <arg> )|\pslabel{verb:clear}} & 
A method with this name will empty or reset the owning object, or it empties a collection that is part of the owning object. It could be used also to \textit{clear} the object indicated by the argument.

Refer to chapter \tqfullref{sec:NamesForCollectionMutators} where \bsq{clear} methods are already mentioned. 

Another sample is \lstinline|Collection::clear|\index{java.util!Collection!clear()} \autocite{ORACLE_DOC_COLLECTION:clear}. \\
\hline

%------------------------------------------------------------------------------

\makecell{\lstinline|clone()|\pslabel{verb:clone}} & 
The method \lstinline|clone()|\index{java.lang!Object!clone()}  \autocite{ORACLE_DOC_OBJECT:clone,ORACLE_DOC_CLONEABLE_INTERFACE} creates an identical copy of the owning object.

Refer to chapter \tqfullvref{sec:Clone} for the details. \\
\hline

%------------------------------------------------------------------------------

\makecell{\lstinline|close()| \\ \lstinline|close…()|\pslabel{verb:close}} & 
Usually, this name will be assigned to methods that close an (external) resource, like a file, a connection or a device. \lstinline|close()|\index{java.lang!AutoCloseable!close()}\index{java.io!Closeable!close()} is the only method of the interfaces \lstinline|Closeable|\index{java.io!Closeable} \autocite{ORACLE_DOC_CLOSEABLE_INTERFACE} and \lstinline|AutoCloseable|\index{java.lang!AutoCloseable} \autocite{ORACLE_DOC_AUTOCLOSEABLE_INTERFACE}.

The latter interface is important for \verb#try-with-resources#\index{try-with-resources} and therefore sometimes the method \lstinline|close()| is used to implement some special features; refer to \tqfullvref{sec:TryWithResources}. \\
\hline

%------------------------------------------------------------------------------

\makecell{\lstinline|compare( <arg> )| \\ \lstinline|compare…( <arg> )|\pslabel{verb:compare}} & 
The method \lstinline|compare()|\index{java.util!Comparator!compare()}  \autocite{ORACLE_DOC_COMPARATOR:compare} is declared by the interface\index{interface} \lstinline|java.util.Comparator|\index{java.util!Comparator} \autocite{ORACLE_DOC_COMPARATOR_INTERFACE}. An implementation of that method compares the given arguments with each other and returns the result as an integer.

Any \bsq{compare} method should be implemented along the lines provided by the documentation for \lstinline|java.util.Comparator.compare()|\index{java.util!Comparator!compare()}, even when the owning class does not implement the interface \lstinline|Comparator|\index{java.util!Comparator}.

See also \hyperref[verb:compareTo]{\bsq{compareTo}}, \hyperref[verb:compareTo]{below}. 

The lambda below illustrates how the interfaces \lstinline|Comparator|\index{java.util!Comparator} and \lstinline|Comparable|\index{java.lang!Comparable} are related to each other:
\index{java.lang!Integer!signum()}
\begin{lstlisting}[xleftmargin=.7cm]
final Comparator<Comparable> comparator = (o1,o2) -> Integer.signum( o1.compareTo( o2 ) );
\end{lstlisting}
\\
\hline

%------------------------------------------------------------------------------

\makecell{\lstinline|compareTo( <arg> )|\pslabel{verb:compareTo}} & 
The method \lstinline|compareTo()|\index{java.lang!Comparable!compareTo()} \autocite{ORACLE_DOC_COMPARABALE:compareTo} is the only method of the interface \lstinline|java.lang.Comparable|\index{java.lang!Comparable} \autocite{ORACLE_DOC_COMPARABALE_INTERFACE}. It compares the owning object with the object provided as the argument and returns the result as an integer, same as a \hyperref[verb:compare]{\bsq{compare}} method. 

You should use the name \lstinline|compareTo()| only for methods that are implementations of the interface \lstinline|Comparable|\index{java.lang!Comparable}. \\
\hline

%------------------------------------------------------------------------------

\makecell{\lstinline|compile( <arg> )| \\ \lstinline|compile…( <arg> )|\pslabel{verb:compile}} & 
A \bsq{compile} method translates a \textit{source} (usually a text) into an executable format and returns that. A sample is \lstinline|Pattern.compile( String )|\index{java.util.regex!Pattern!compile()} \autocite{ORACLE_DOC_PATTERN:compile}. \\
\hline

%------------------------------------------------------------------------------

\makecell{\lstinline|compose…( <arg> )| \\ \lstinline|compose…()|\pslabel{verb:compose}} & 
A \bsq{compose} method will compose a new object or data structure from the internal state of the owning object and the given arguments – if any. It will not modify the owning object itself, and each call with identical arguments (and on an unchanged object) will return the same value.

Other than a \hyperref[verb:build]{\bsq{build}} method, a \bsq{compose} method can be called more than once on the same object instance. \\
\hline

%------------------------------------------------------------------------------

\makecell{\lstinline|connect()| \\ \lstinline|connect( <arg> )| \\ \lstinline|connect…()| \\ \lstinline|connect…( <arg> )|\pslabel{verb:connect}} & 
A method with this name initialises a connection from the owning object to an external resource, like a device or a remote system. \\
\hline

%------------------------------------------------------------------------------

\makecell{\lstinline|convert…()| \\ \lstinline|convert…( <arg> )| \\ \lstinline|convert( <arg> )|\pslabel{verb:convert}} & 
A convert method transforms something from one format into another. The details are determined by the arguments and the state of the owning object.

The name is often supplemented by \bsq{from} or \bsq{to}, to indicate the source or target format, as in \lstinline|convertFromXML()| or \lstinline|convertToHTML()|. \\
\hline

%------------------------------------------------------------------------------

\makecell{\lstinline|copy( <arg> )| \\ \lstinline|copyOf( <arg> )| \\ \lstinline|copy…( <arg> )|\pslabel{verb:copy}} & 
A \bsq{copy} method makes a semanticall copy of the given argument and returns it. If this is also an identical copy (like through the invocation of \hyperref[verb:clone]{\lstinline|clone()|}\index{java.lang!Object!clone()}) or not depends from the implementation. 

One example is \lstinline|List::copyOf|\index{java.util!List!copyOf()} \autocite{ORACLE_DOC_LIST:copyOf}, another example is \lstinline|Collections.copy(List,List)|\index{java.util!Collections!copy()} \autocite{ORACLE_DOC_COLLECTIONS:copy}. \\
\hline

%------------------------------------------------------------------------------

\makecell{\lstinline|create…( <arg> )| \\ \lstinline|create…()|\pslabel{verb:create}} & 
This name indicates a method that creates something based on the given arguments and/or on the state of the owning object, but different from the \hyperref[verb:compose]{\bsq{compose}} method, each call can have a different result, even for identical arguments.

A \bsq{create} method usually returns the created entity as the result. \\
\hline

%------------------------------------------------------------------------------

\makecell{\lstinline|decrement…()|/\lstinline|dec…()| \\ \lstinline|decrement()|/\lstinline|dec()|\pslabel{verb:decrement}} & 
A method with the name \lstinline|decrement()| reduces the value of the owning object by 1. If \bsq{decrement} is used as prefix to the name, the name indicates the object property whose value is reduced by 1.

The converse operation is named \hyperref[verb:increment]{\bsq{increment}}. \\
\hline

%------------------------------------------------------------------------------

\makecell{\lstinline|delete()| \\ \lstinline|delete( <arg> )| \\ \lstinline|delete…()| \\ \lstinline|delete…( <arg> )|\pslabel{verb:delete}} & 
A method with the name \bsq{delete} deletes an external entity (in most cases, a file or folder), either identified by the state of the owning object or by the given arguments.

An example is \lstinline|java.io.File.delete()|\index{java.io!File!delete()} \autocite{ORACLE_DOC_FILE:delete} that removes the file that is associated with the respective instance of \lstinline|java.io.File|\index{java.io!File} \autocite{ORACLE_DOC_FILE_CLASS}. 

Another example is \lstinline|java.nio.file.Files::delete|\index{java.nio.file!Files!delete()} \autocite{ORACLE_DOC_FILES:delete} that deletes the file that is specified by the given \lstinline|Path|\index{java.nio.file.Path} \autocite{ORACLE_DOC_PATH_INTERFACE} argument. \\
\hline

%------------------------------------------------------------------------------

\makecell{\lstinline|draw…( <arg> )| \\ \lstinline|draw( <arg> )|\pslabel{verb:draw}} & 
The prefixes \bsq{draw} and \hyperref[verb:paint]{\bsq{paint}} are typically found for method names related to elements of graphical user interfaces\index{graphical user interface|see {GUI}} (GUI\index{GUI}).

Both are part of the names for methods that will visualise GUI\index{GUI} elements on screen; the difference is that the \bsq{draw} methods are usually associated with the container and take the element to draw as the argument, while the \bsq{paint} methods are associated with the element and take the container where to paint that element. 

Examples are 
\begin{itemize}[nosep]
\item{\lstinline|java.awt.Graphics2D::draw|\index{java.awt!Graphics2D!draw()} \autocite{ORACLE_DOC_GRAPHICS2D:draw}}
\item{\lstinline|javax.swing.text.PlainView::drawLine|\index{javax.swing.text!PlainView!drawLine} \autocite{ORACLE_DOC_PLAINVIEW:drawLine}}
\item{\lstinline|javax.scene.canvas.GraphicsContext::drawImage|\index{javafx.scene.canvas!GraphicsContext!drawImage()} \autocite{JavaFX_DOC_GraphicsContext:drawImage}}
\item{\lstinline|java.awt.Window::paint|\index{java.awt!Window!paint()} \autocite{ORACLE_DOC_WINDOW:paint}}
\item{\lstinline|javax.swing.JComponent:paint|\index{javax.swing!JComponent!paint()} \autocite{ORACLE_DOC_JCOMPONENT:paint}}
\item{\lstinline|javafx.embed.swing.JFXPanel::paintComponent|\index{javax.embed.swing!JFXPanel!paintComponent()} \autocite{JavaFX_DOC_JFXPanel:paintComponent}}
\end{itemize} \\
\hline

%------------------------------------------------------------------------------

\makecell{\lstinline|exec…()| \\ \lstinline|exec…( <arg> )|\pslabel{verb:exec}} & 
Use the prefix \bsq{exec} to name methods that will execute an operation, identified by the main part of the name. The arguments are parameters for this operation.

Nothing is said about the state of the owning object, any return values or the state of the arguments.

See also \hyperref[verb:action]{\bsq{action}} and \hyperref[verb:perform]{\bsq{perform}}. \\
\hline

%------------------------------------------------------------------------------

\makecell{\lstinline|execute()| \\ \lstinline|execute( <arg> )|\pslabel{verb:execute}} & 
This is the name for a method that does the program's or application's work. The arguments are usually the command line arguments for the program, although this is not mandatory.

This is used in a pattern similar to this:
\begin{lstlisting}[xleftmargin=.7cm]
public final class MyProgram
{
    public final void execute( final String [] args ) { … }
    public final boolean initialize( final String [] args ) { … }
    public static final void main( final String... args )
    {
        final var application = new MyProgram();
        if( application.initialize( args ) ) application.execute();
    }   //  main()
}
//  class MyProgram
\end{lstlisting}

The method \hyperref[verb:initialize]{\lstinline|initialize()|} is optional, but if available, it takes the command line arguments instead of \lstinline|execute()|. \lstinline|execute()| may take arguments that where generated by \lstinline|initialize()|. If the method is public, this would allow to use the respective class not only as a standalone program, but also in the context of another program/application. \\
\hline

%------------------------------------------------------------------------------

\makecell{\lstinline|exit( <arg> )|\pslabel{verb:exit}} & 
The class \lstinline|System|\index{java.lang!System} \autocite{ORACLE_DOC_SYSTEM_CLASS} defines the method \lstinline|exit(int)|\index{java.lang!System!exit()} \autocite{ORACLE_DOC_SYSTEM:exit} that, when called, terminates the JVM. Depending on the type of product\footnote{refer to \tqfullvref{sec:TypesOfProducts}.} that your project delivers, this method is seen as evil, just bad or even useful.

The method \lstinline|javafx.application.Platform.exit()|\index{javafx.application!Platform!exit()} \autocite{JavaFX_DOC_Platform:exit} has basically the same effect.

You should avoid naming a method \lstinline|exit()|, and you should even avoid to use \bsq{exit} as a prefix of a method name, just to avoid a misunderstanding.

But naming a method \lstinline|exit()| that just does some housekeeping before it will call \lstinline|System::exit|\index{java.lang!System!exit()} is perfectly ok. \\
\hline

%------------------------------------------------------------------------------

\makecell{\lstinline|filter…( <arg> )| \\ \lstinline|filter( <arg> )|\pslabel{verb:filter}} & 
A \bsq{filter} method will be applied on a collection and returns those that matches a filter condition. The collection or the filter condition or both are provided as arguments. 

An example is the method \lstinline|Stream::filter|\index{java.util.stream!Stream!filter()} \autocite{ORACLE_DOC_STREAM:filter}. It works on an instance of \lstinline|java.util.stream.Stream|\index{java.util.stream!Stream} \autocite{ORACLE_DOC_STREAM_INTERFACE}, takes a \lstinline|Predicate|\index{java.util.function!Predicate} as the filter condition, and returns another \lstinline|Stream| with those values that fullfilled the condition. 

Other \bsq{filter} methods should basically work along the same lines. \\
\hline

%------------------------------------------------------------------------------

\makecell{\lstinline|find( <arg> )| \\ \lstinline|find…( <arg> )| \\ \lstinline|find…()|\pslabel{verb:find}} & 
A \bsq{find} method searches something that is expected to be there (different from a \hyperref[verb:search]{\bsq{search}} method described below). The method returns the result, or it throws an appropriate exception when it fails. It would be also possible to return an instance of \lstinline|java.lang.Optional|\index{java.lang!Optional} \autocite{ORACLE_DOC_OPTIONAL_CLASS} with the result.

The arguments are parameters for the search, although it is possible that the state of the owning object determines the search.

The method may or may not change the state of the owning object. \\
\hline

%------------------------------------------------------------------------------

\makecell{\lstinline|format…( <arg> )| \\ \lstinline|format( <arg> )|\pslabel{verb:format}} & 
A \bsq{format} method prepares the representation of data. A wonderful example is the method \lstinline|format|\index{java.util!Formatter!format()} \autocite{ORACLE_DOC_FORMATTER:format} from the class \lstinline|java.util.Formatter|\index{java.util!Formatter} \autocite{ORACLE_DOC_FORMATTER_CLASS} (the method \lstinline|String::format|\index{java.lang!String!format} \autocite{ORACLE_DOC_STRING:format} delegates to this), another one is \lstinline|format()|\index{java.time.format!DateTimeFormatter!format()} \autocite{ORACLE_DOC_DATETIMEFORMATTER:format} from class \lstinline|java.time.format.DateTimeFormatter|\index{java.time.format!DateTimeFormatter} \autocite{ORACLE_DOC_DATETIMEFORMATTER_CLASS}. \\
\hline

%------------------------------------------------------------------------------

\makecell{\lstinline|from( <arg> )| \\ \lstinline|from…( <arg> )| \\ \lstinline|...From( <arg> )|\pslabel{verb:from}} & 
Methods with the name \lstinline|from()| or a name prefixed with \bsq{from} are (usually static) factory methods\index{method!factory method}\index{factory} that create an instance of the owning class from something, represented by the given argument.

A sample would be \lstinline|Instant.from()|\index{java.time!Instant!from()} \autocite{ORACLE_DOC_INSTANT:from}.

\bsq{From} can be found also a suffix for method names, as in \lstinline|Signature.parseFrom()|\index{java.lang.classfile!Signature!parseFrom()} \autocite{ORACLE_DOC_SIGNATURE:parseFrom} (although this would be a better sample for a \hyperref[verb:parse]{\bsq{parse}} method). \\
\hline

%------------------------------------------------------------------------------

\makecell{\lstinline|generate…()| \\ \lstinline|generate…( <arg> )| \\ \lstinline|generate()| \\ \lstinline|generate( <arg> )|\pslabel{verb:generate}} & 
A \bsq{generate} method generates something based on the state of the owning object and the given arguments, with or without changing that state. Each call to a \bsq{generate} method may have a different result, even for the same arguments – this distinguishes it from \hyperref[verb:compose]{\bsq{compose}} methods.

Different from a \hyperref[verb:create]{\bsq{create}} method, a \bsq{generate} method may work on entities outside the program, like a file, an image, a report, or alike. This means it does not necessarily have a return value, or the return value may just be the status of the generation process. \\
\hline

%------------------------------------------------------------------------------

\makecell{\lstinline|get…()|\pslabel{verb:get}} & 
The prefix \bsq{get} indicates a \textit{getter} method\index{method!getter method} if the respective method does not take any arguments. It does not change the state of the owning object.

The remaining part of the method name is known as the \bsq{property name}\index{property}\footnote{refer to chapter \tqfullvref{sec:NamesOfProperties}} the getter\index{method!getter method} returns the value for that property (attribute). For more details, refer to \autocite{ORACLE_DOC_JAVABEANS:Chapter8_3}\footnote{Although the whole document \autocite{ORACLE_DOC_JAVABEANS} could be relevant here.}

For optional properties\index{property}, a getter\index{method!getter method} may return an instance of \lstinline|java.util.Optional|\index{java.util!Optional} \autocite{ORACLE_DOC_OPTIONAL_CLASS}, but it should never fail (meaning it should be avoided that a getter throws an exception). \\
\hline

%------------------------------------------------------------------------------

\makecell{\lstinline|get…( <arg> )| \\ \lstinline|get( <arg> )|\pslabel{verb:get_Arg}} & 
A \bsq{get} method that takes an argument is \textit{not} a getter method\index{method!getter method}; nevertheless, it should also not change the state of the owning object.

Such a method returns something that is determined by the rest of the name and further specified by the arguments.

An example for this kind of \bsq{get} is \lstinline|map.get()|\index{java.util!Map!get()} \autocite{ORACLE_DOC_MAP:get}, more samples can be found in \lstinline|ResultSet|\index{java.sql!ResultSet} \autocite{ORACLE_DOC_RESULTSET_INTERFACE}. \\
\hline

%------------------------------------------------------------------------------

\makecell{\lstinline|getInstance()|\pslabel{verb:getInstance}} & 
The method \lstinline|public static <Type> getInstance()| is somehow special; it is not really a getter\index{method!getter method} although it has the prefix \bsq{get} and does not take any arguments.

The method returns an instance of the class \lstinline|<Type>|; each invocation returns the same instance but never \lstinline|null|. It does not matter whether the first invocation creates that instance, or if it will exist previously. 

If declared by an interface\index{interface} or an abstract class\index{class!abstract class}, usually it will be a factory method\index{factory}\index{method!factory method} for an object instance that implements that interface\index{interface} or abstract class\index{class!abstract class}. A factory\index{factory}\index{method!factory method} that returns a new instance for each invocation should be named \hyperref[verb:newInstance]{\lstinline|newInstance()|} instead. \\
\hline

%------------------------------------------------------------------------------

\makecell{\lstinline|has…()| \\ \lstinline|has( <arg> )| \\ \lstinline|has…( <arg> )|\pslabel{verb:has}} & 
A method with a name starting with \bsq{has} is meant to check whether the owning object \bdq{contains} something, either identified by the rest of the name, or by the given argument. Samples could be \lstinline|hasChildren()| for a tree node object, or \lstinline|hasOrders()| for a customer record instance.

Calling the method may not change the state of the owning object. \\
\hline

%------------------------------------------------------------------------------

\makecell{\lstinline|hasNext()|\pslabel{verb:hasNext}} & 
The method \lstinline|hasNext|\index{java.util!Iterator!hasNext()} \autocite{ORACLE_DOC_ITERATOR:hasNext} is declared by the interface \lstinline|java.util.Iterator|\index{java.util!Iterator} \autocite{ORACLE_DOC_ITERATOR_INTERFACE}. It checks whether the iterator can still provide at least one more value. 

A method with that name in a class that does not implement \lstinline|Iterator|\index{java.util!Iterator} should do something similar. \\
\hline

%------------------------------------------------------------------------------

\makecell{\lstinline|increment…()|/\lstinline|inc…()| \\ \lstinline|increment()|/\lstinline|inc()|\pslabel{verb:increment}} & 
A method with the name \lstinline|increment()| increases the value of the owning object by 1. If \bsq{increment} is used as prefix to the name, the name indicates the object property whose value is increased by 1.

The converse operation is named \hyperref[verb:decrement]{\bsq{decrement}}. \\
\hline

%------------------------------------------------------------------------------

\makecell{\lstinline|initialize()| \\ \lstinline|initialize( <arg> )|\pslabel{verb:initialize}} & 
The \bsq{initialize} method does what its name says: it initialises the owning object. Obviously, this changes the state of that object.

If the method cannot be called multiple times, it should throw an \lstinline|IllegalStateException|\index{java.lang!IllegalStateException} \autocite{ORACLE_DOC_ILLEGALSTATEEXCEPTION_CLASS} on any invocation after the first. \\
\hline

%------------------------------------------------------------------------------

\makecell{\lstinline|insert…( <arg> )| \\ \lstinline|insert( <arg> )|\pslabel{verb:insert}} & 
An \bsq{insert} method implements data within a previously existing sequence of data. The SQL command \lstinline|INSERT| is a good example of that, but with the method \lstinline|insert()|\index{java.lang!StringBuilder!insert()} \autocite{ORACLE_DOC_STRINGBUILDER:insert} from the class \lstinline|java.lang.StringBuilder|\index{java.lang!StringBuilder} \autocite{ORACLE_DOC_STRINGBUILDER_CLASS}, there is also a good sample from the Java runtime library. \\
\hline

%------------------------------------------------------------------------------

\makecell{\lstinline|is…()| \\ \lstinline|is…( <arg> )|\pslabel{verb:is}} & 
Usually, \bsq{is} is the prefix for a \textit{getter}\index{method!getter method} that returns a \lstinline|boolean| (refer to \hyperref[verb:get]{\bsq{get}}, above), but it can be used also to query the state of the owning object, or the state of a given argument, like in the samples below:
\index{java.time!DayOfWeek}
\begin{lstlisting}[xleftmargin=.7cm]
public static final boolean isWorkday( final DayOfWeek day ) { return day != DayOfWeek.SUNDAY; } 
public static final boolean isWeekend( final DayOfWeek day ) { return day == DayOfWeek.SATURDAY || day == DayOfWeek.SUNDAY; } 
\end{lstlisting} \\
    \hline

%------------------------------------------------------------------------------

\makecell{\lstinline|join()| \\ \lstinline|join( <arg> )|\pslabel{verb:join}} & 
The method \lstinline|Thread.join()|\index{java.lang!Thread!join()} \autocite{ORACLE_DOC_THREAD:join} pauses the calling thread (the current thread) until the thread, represented by the \lstinline|Thread|\index{java.lang!Thread} \autocite{ORACLE_DOC_THREAD_CLASS} object that owns the methods, terminates, or, in case of the variant that takes an argument, a timeout occurs.

If you implement a \bsq{join} method in one of your own classes, it should do something similar, to avoid confusion. 

Unfortunately, the Java runtime library has several classes that implement a \bsq{join} method with different functionality. Most prominent is \lstinline|String::join|\index{java.lang!String!join()} \autocite{ORACLE_DOC_STRING:join} that is used to compose a single \lstinline|String|\index{java.lang!String} from a list of \lstinline|CharSequence|\index{java.lang!CharSequence} instances. \\
\hline

%------------------------------------------------------------------------------

\makecell{\lstinline|load()| \\ \lstinline|load( <arg> )|\pslabel{verb:load}} & 
A method with the name \lstinline|load()| \bdq{loads} data from somewhere (usually an external source) and initialises the owning object with this data. The argument is usually the reference to the source.

The return type is either \lstinline|void|, or the method returns some kind of status for the load operation. \\
\hline

%------------------------------------------------------------------------------

\makecell{\lstinline|load…()| \\ \lstinline|load…( <arg> )|\pslabel{verb:load_name}} & 
A method with a name prefixed by \bsq{load} also loads data from an external source (like \hyperref[verb:load]{\bsq{load} above}), but it initialises and returns a data structure that is described by the remainder of the method name.

Depending on the context, a name like \lstinline|loadData()| could be valid, but some more meaningful is preferred, like \lstinline|loadClientHistory()| or \lstinline|loadArticleList()|. \\
\hline

%------------------------------------------------------------------------------

\makecell{\lstinline|log( <arg> )| \\ \lstinline|log…( <arg> )| \\ \lstinline|log()|\pslabel{verb:log}} & 
A \bsq{log} method writes something to a logging subsystem of some kind. The method is either a member of the class representing that subsystem, or it is a member of an arbitrary class. In the first case, the method's arguments determine what is logged. In the latter case, the method logs the internal state of the owning object, without modifying it. \\
\hline

%------------------------------------------------------------------------------

\makecell{\lstinline|main()| \\ \lstinline|main( <arg> )|\pslabel{verb:main}} & 
The method \lstinline|main()| is the entry point for any Java program; refer to \autocite{ORACLE_DOC_LANGUAGE_SPECIFICATION:Invoke_main} for the details about the \lstinline|main()| method.

The signature for the \lstinline|main()| method that is found most often in literature about the \Java programming language is
\begin{lstlisting}[xleftmargin=.7cm]
public static void main( String [] args )
\end{lstlisting}
but I prefer it as
\begin{lstlisting}[xleftmargin=.7cm]
public static final void main( final String... args )
\end{lstlisting}
because it reflects the command line better. Refer to \autocite{ORACLE_DOC_LANGUAGE_SPECIFICATION:Invoke_Test.main} about the signatures. \\
\hline

%------------------------------------------------------------------------------

\makecell{\lstinline|make…()| \\ \lstinline|make…( <arg> )|\pslabel{verb:make}} & 
The prefix \bsq{make} is similar to \hyperref[verb:create]{\bsq{create}}, \hyperref[verb:compose]{\bsq{compose}}, or \hyperref[verb:generate]{\bsq{generate}}, but usually a method with a name starting with \bsq{make} produces something outside the program – like a file or a database (table). In this regard it is like \bsq{generate}, but the process of \bdq{making} is less complex as that of \bdq{generation}. \\
\hline

%------------------------------------------------------------------------------

\makecell{\lstinline|mark()| \\ \lstinline|mark( <arg> )| \\ \lstinline|mark…()| \\ \lstinline|mark…( <arg> )|\pslabel{verb:mark}} & 
A \bsq{mark} method changes the state of the owning object by setting (when taking a \lstinline|boolean| argument) or toggling (when taking no argument) an internal flag (see also \hyperref[verb:check]{\bsq{check}}). 

Usually those methods will not return something, but for chaining, they can return \lstinline|this|, or they return the previous state of the flag. \\
\hline

%------------------------------------------------------------------------------

\makecell{\lstinline|map…( <arg> )| \\ \lstinline|map( <arg> )|\pslabel{verb:map}} & 
A \bsq{map} method translates something into something else, like the method \lstinline|Optional::map|\index{java.util!Optional!map()} \autocite{ORACLE_DOC_OPTIONAL:map}.

Another example would be the method \lstinline|mapFromNull()| below:
\begin{lstlisting}[caption=mapFromNull(),xleftmargin=.7cm]
public static final <T> T mapFromNull( final T value, final Supplier<T> supplier )
{
    return isNull( value ) ? supplier.get() : value;
}   //  mapFromNull() 
\end{lstlisting} \\
\hline

%------------------------------------------------------------------------------

\makecell{\lstinline|merge…( <arg> )| \\ \lstinline|merge( <arg> )|\pslabel{verb:merge}} & 
Merge methods are combining two sets of data into one single one. If both data sets are provided as argument, a merge method usually does not change the state of the owning object, or it can even be static.

Examples are \lstinline|StringJoiner::merge|\index{java.util!StringJoiner!merge()} \autocite{ORACLE_DOC_STRINGJOINER:merge} and \lstinline|Map::merge|\index{java.util!Map!merge()} \autocite{ORACLE_DOC_MAP:merge}. \\
\hline

%------------------------------------------------------------------------------

\makecell{\lstinline|new…( <arg> )| \\ \lstinline|new…()|\pslabel{verb:new}} & 
A method \lstinline|new()|, with or without arguments, is invalid because \lstinline|new| is a keyword of the \Java Programming Language (refer to \autocite{ORACLE_DOC_LANGUAGE_SPECIFICATION:Keywords}). The method reference\lstinline|<Class>::new| refers to the constructor\index{constructor} of \lstinline|class| that has the appropriate parameters as inferred from the context. This might look like line~3:
\index{java.util!Set}\index{java.util!HashSet}\index{java.util!Set!toArray()}
\begin{lstlisting}[xleftmargin=.7cm]
final Set<String> strings = new HashSet<>();
…
final var stringArray = strings.toArray( String []::new );
\end{lstlisting}
The method \lstinline|Set::toArray|\index{java.util!Set!toArray()} is declared as below (inherited from \lstinline|Collection|\index{java.util!Collection} \autocite{ORACLE_DOC_COLLECTION:toArray3}):
\begin{lstlisting}[xleftmargin=.7cm]
public <T> T[] toArray( IntFunction<T[]> generator );
\end{lstlisting}

Consequently, \bsq{new} can be used only as the prefix for a method name, and that should be a factory method\index{method!factory method} for something that is determined by the remainder of the method name.

See also \hyperref[verb:newInstance]{below}. \\
\hline

%------------------------------------------------------------------------------

\makecell{\lstinline|newInstance( <arg> )| \\ \lstinline|newInstance()|\pslabel{verb:newInstance}} & 
The method \lstinline|newInstance()| is usually a static member of a class (probably abstract) or an interface, and it returns a new instance of the owning class when called (it is a factory method\index{method!factory method} for that class or interface). Depending on the arguments for that call, the instances might be equals, but they will never be the same; this distinguishes \lstinline|newInstance()| from \hyperref[verb:getInstance]{\lstinline|getInstance()|}. \\
\hline

%------------------------------------------------------------------------------

\makecell{\lstinline|obtain…()| \\ \lstinline|obtain…( <arg> )|\pslabel{verb:obtain}} & 
A method with a name prefixed with \bsq{obtain} is similar to a \hyperref[verb:get]{getter}\index{method!getter method}, but with the difference that an \bsq{obtain} method can change the state of the owning object, and that it does not necessarily return just a property\index{property}.

An \bsq{obtain} method should never fail (should not throw an exception), but it may return an instance of \lstinline|java.util.Optional|\autocite{ORACLE_DOC_OPTIONAL_CLASS}. \\
\hline

%------------------------------------------------------------------------------

\makecell{\lstinline|of( <arg> )| \\ \lstinline|of…( <arg> )| \\ \lstinline|...Of( <arg> )|\pslabel{verb:of}} & 
The \bsq{of} methods are quite similar to the \hyperref[verb:from]{\bsq{from}} methods, and like those or \hyperref[verb:newInstance]{\lstinline|newInstance()|}, these methods create an instance of the owning class or interface based on the given arguments. \\
\hline

%------------------------------------------------------------------------------

\makecell{\lstinline|open()|\pslabel{verb:open}} & 
A method \lstinline|open| opens a connection to an external resource represented by the owning object of that method. Obviously, the object's state will change when the method is called.

Usually, \lstinline|open()| returns nothing, or a status of the open operation.

The method will also not take any arguments as all necessary information are given by the state of the owning object. Usually the defining class will provide also a corresponding \hyperref[verb:close]{\lstinline|close()|} method. \\
\hline

%------------------------------------------------------------------------------

\makecell{\lstinline|open…( <arg> )| \\ \lstinline|open( <arg> )|\pslabel{verb:opern_arg}} & 
A method with a name that is prefixed with \bsq{open} will also open a connection to an external resource, but it will return an object representing either that connection or the external resource. The required information can be taken from the owning object or given as the arguments or both.

Such a method may or may not change the state of the owning object.

The corresponding \hyperref[verb:close]{\bsq{close}} method is usually provided by the returned object.

A sample could be a method
\begin{lstlisting}
public final File openLogFile( final File logFolder ) {…}
\end{lstlisting}
 \\
\hline

%------------------------------------------------------------------------------

\makecell{\lstinline|paint…( <arg> )| \\ \lstinline|paint( <arg> )|\pslabel{verb:paint}} & 
Refer to \hyperref[verb:draw]{\bsq{draw}}. \\
\hline

%------------------------------------------------------------------------------

\makecell{\lstinline|peek()|\pslabel{verb:peek}} & 
A method with this name is meant to spy into a data structure without modifying the state of the owning object. The method \lstinline|java.util.Queue::peek|\index{java.util!Queue!peek()} \autocite{ORACLE_DOC_QUEUE:peek} is a sample for this.

The method \lstinline|java.util.stream.Stream::peek|\index{java.util.stream!Stream!peek()} \autocite{ORACLE_DOC_STREAM:peek} does not examine a data structure, but it allows to spy on the intermediate results of the stream processing. \\
\hline

%------------------------------------------------------------------------------

\makecell{\lstinline|perform()| \\ \lstinline|perform( <arg> )| \\ \lstinline|perform()| \\ \lstinline|perform( <arg> )|\pslabel{verb:perform}} & 
A \bsq{perform} method performs an action, either implicitly like for an implementation of the command pattern, or explicitly named by the remaining part of the method name; it is even an option that the action is given as the argument of the method. Usually such a method will not change the state if the owning object.

Refer also to \hyperref[verb:action]{\bsq{action}} and \hyperref[verb:exec]{\bsq{exec}}. \\
\hline

%------------------------------------------------------------------------------

\makecell{\lstinline|pop()| \\ \lstinline|pop…()|\pslabel{verb:pop}} & 
For the data structure \bdq{Stack}, \bsq{pop} is the operation that takes the top entry from the stack. The converse operation is \hyperref[verb:push]{\bsq{push}}. 

Methods with this name or \bsq{pop} as prefix for their name should do something similar. \\
\hline

%------------------------------------------------------------------------------

\makecell{\lstinline|print( <arg> )| \\ \lstinline|print()|\pslabel{verb:print}} & 
A print method will write the textual representation of either the argument to the output represented by the owning object, or of the owning object to the output represented by the argument.

An example is \lstinline|PrintStream::print|\index{java.io!PrintStream!print()} \autocite{ORACLE_DOC_PRINTSTREAM:print}; the IOStream classes and interfaces define a whole family of \bsq{print} methods. \\
\hline

%------------------------------------------------------------------------------

\makecell{\lstinline|pull()| \\ \lstinline|pull( <arg> )| \\ \lstinline|pull…()| \\ \lstinline|pull…( <arg> )|\pslabel{verb:pull}} & 
A \bsq{pull} method retrieves data from somewhere, most probably from an external resource. It is somehow related to \hyperref[verb:obtain]{\bsq{obtain}} or \hyperref[verb:retrieve]{\bsq{retrieve}}.
 \\
\hline

%------------------------------------------------------------------------------

\makecell{\lstinline|push( <arg> )| \\ \lstinline|push…( <arg> )|\pslabel{verb:push}} & 
For the data structure \bdq{Stack}, \bsq{push} is the operation that puts a new entry to the top of the stack. The converse operation is \hyperref[verb:pop]{\bsq{pop}}. 

Methods with this name or \bsq{push} as prefix for their name should do something similar. \\
\hline

%------------------------------------------------------------------------------

\makecell{\lstinline|put( <arg> )| \\ \lstinline|put…( <arg> )|\pslabel{verb:put}} & 
\bsq{put} is somehow a synonym for \hyperref[verb:add]{\bsq{add}}, but different from that it has some position identification: with \hyperref[verb:add]{\bsq{add}}, the new part is just thrown on the big pile, while it will be placed to a special position with \bsq{put}. Refer to the method \lstinline|map::put|\index{java.util!Map!put()} \autocite{ORACLE_DOC_MAP:put}. \\
\hline

%------------------------------------------------------------------------------

    \makecell{\lstinline|read()|\pslabel{verb:read}} & 
    The \lstinline|read()| method loads the state of the owning object from an external resource. Usually, it does not return something, or it returns that status of the ‘read’ operation.\newline I prefer to use \lstinline|load()| instead. \\
    \hline

    \makecell{\lstinline|read( <arg> )| \\ \lstinline|read…( <arg> )|\pslabel{verb:read_arg}} & 
    A ‘read’ method read something from an external source and returns. This may or may not change the state of the owning object.\newline Usually a reference to that external source is provided as a parameter, but it is also possible that this reference is implicit (as for streams). \\
    \hline

%------------------------------------------------------------------------------

\makecell{\lstinline|reduce…( <arg> )| \\ \lstinline|reduce( <arg> )|\pslabel{verb:reduce}} & 
\lipsum[1] \\
\hline

%------------------------------------------------------------------------------

    \makecell{\lstinline|release…( <arg> )| \\ \lstinline|release( <arg> )|\pslabel{verb:release}} & 
A ‘release’ method provides the converse operation to an \hyperref[verb:acquire]{\bsq{acquire}} method: it releases the exclusive access to the object that is somehow identified by the given argument. It does not return something, or it returns the status of the operation. \\
    \hline

    \makecell{\lstinline|remove…( <arg> )| \\ \lstinline|remove( <arg> )|\pslabel{verb:remove}} & 
    A method with the name \lstinline|remove()| or a name that is prefixed with ‘remove’ will remove something from an internal data structure of the owning method. The given argument is some kind of reference to that thing that should be removed. It does not return something, or it returns the removed object or the status of the operation.\newline A ‘remove’ method is somehow the reverse operation to an ‘add’ method. \\
    \hline

%------------------------------------------------------------------------------

\makecell{\lstinline|repose…( <arg> )| \\ \lstinline|repose( <arg> )|\pslabel{verb:repose}} & 
\lipsum[1] \\
\hline

%------------------------------------------------------------------------------

\makecell{\lstinline|reset…( <arg> )| \\ \lstinline|reset( <arg> )|\pslabel{verb:reset}} & 
\lipsum[1] \\
\hline

%------------------------------------------------------------------------------

    \makecell{\lstinline|retrieve…()| \\ \lstinline|retrieve…( <arg> )|\pslabel{verb:retrieve}} & 
    A ‘retrieve’ method is similar to a ‘get’ method, with the differences that it may change the state of the owning object, that it is not limited in returning a property, and – that makes it also different to an ‘obtain’ method – that it can fail. \\
    \hline

    \makecell{\lstinline|run()|\pslabel{verb:run}} & \lstinline|run()| is the name of the main method of all implementations of the interface \lstinline|java.lang.Runnable|\autocite{ORACLE_DOC_RUNNABLE_INTERFACE}, and as such it is also the main method of all threads except the main thread (there it is \lstinline|main()|). The name should not be used in a different context. Refer also to \autocite{ORACLE_DOC_RUNNABLE:run,ORACLE_DOC_THREAD:run}. \\
    \hline

%------------------------------------------------------------------------------

\makecell{\lstinline|search( <arg> )| \\ \lstinline|search…( <arg> )|\pslabel{verb:search}} & 
A \bsq{search} method searches something, but it is a valid result that it does not find something (it will never throw an exception if nothing was found). That distinguishes it from a \hyperref[verb:find]{\bsq{find}} method as described above. The method returns an instance of \lstinline|java.util.Optional|\index{java.util!Optional} \autocite{ORACLE_DOC_OPTIONAL_CLASS} with the result.

The arguments are parameters for the search.

The method may or may not change the state of the owning object. \\
    \hline

%------------------------------------------------------------------------------

    \makecell{\lstinline|save()| \\ \lstinline|save( <arg> )|\pslabel{verb:save}} & 
    A method \lstinline|save()| provides the converse operation to that of a method \lstinline|load()|: it persists the current state of the owning object, usually to an external destination. The argument usually determines that destination.\newline If the method has another return type than \lstinline|void|, the return value is the status of the save operation.\newline The method may or may not change the state of the owning object. \\
    \hline

    \makecell{\lstinline|save…( <arg> )|\pslabel{verb:save_arg}} & 
    A ‘save’ method provides the converse operation to that of a ‘load’ method: it persists the object that is identified by the given argument to an external destination. If it has a return type other than \lstinline|void|, this will be the status of the save operation.\newline The method may or may not change the state of the owning object. \\
    \hline

%------------------------------------------------------------------------------

\makecell{\lstinline|select…( <arg> )| \\ \lstinline|select( <arg> )|\pslabel{verb:select}} & 
\lipsum[1] \\
\hline

%------------------------------------------------------------------------------

\makecell{\lstinline|shutdown…( <arg> )| \\ \lstinline|shutdown( <arg> )|\pslabel{verb:shutdown}} & 
\lipsum[1] \\
\hline

%------------------------------------------------------------------------------

\makecell{\lstinline|sleep…( <arg> )| \\ \lstinline|sleep( <arg> )|\pslabel{verb:sleep}} & 
\lipsum[1] \\
\hline

%------------------------------------------------------------------------------

\makecell{\lstinline|split…( <arg> )| \\ \lstinline|split( <arg> )|\pslabel{verb:split}} & 
\lipsum[1] \\
\hline

%------------------------------------------------------------------------------

\makecell{\lstinline|switch…( <arg> )| \\ \lstinline|switch…()|\pslabel{verb:switch}} & 
A method \lstinline|switch()|, with or without arguments, is invalid because \lstinline|switch| is a keyword of the \Java Programming Language (refer to \autocite{ORACLE_DOC_LANGUAGE_SPECIFICATION:Keywords}). Consequently, \bsq{switch} can be used only a the prefix for a method name.

Use \bsq{switch} for methods that somehow alter the behaviour of something, like \lstinline|switchDirection()| or \lstinline|switchContext()|. \\
\hline

%------------------------------------------------------------------------------

\makecell{\lstinline|terminate…( <arg> )| \\ \lstinline|terminate( <arg> )|\pslabel{verb:terminate}} & 
\lipsum[1] \\
\hline

%------------------------------------------------------------------------------

\makecell{\lstinline|toggle…( <arg> )| \\ \lstinline|toggle( <arg> )|\pslabel{verb:toggle}} & 
\lipsum[1] \\
\hline

%------------------------------------------------------------------------------

\makecell{\lstinline|update…( <arg> )| \\ \lstinline|update( <arg> )|\pslabel{verb:update}} & 
\lipsum[1] \\
\hline

%------------------------------------------------------------------------------

\makecell{\lstinline|wait…( <arg> )| \\ \lstinline|wait( <arg> )|\pslabel{verb:wait}} & 
\lipsum[1] \\
\hline

%------------------------------------------------------------------------------

\end{longtable}

%==============================================================================
% End of file!!
%==============================================================================
